\chapter{DF-062 Audit notification system family across real surfaces}\label{record-103}

\noindent{\footnotesize\color{BookMuted}\path{docs/design/DF-062-audit-notification-system-family-across-real-surfaces.md}\par}\smallskip

\begin{RecordMeta}
title: DF-062 Audit notification system family across real surfaces
legend: DF
lane: design
priority: medium
keywords:
  - component-audit
  - dogfood
  - lowerings
  - notifications

\end{RecordMeta}

\section{Framing}\label{record-103-framing}

The notification system is the family for transient events that need more structure than one directly composed \texttt{t\allowbreak{}o\allowbreak{}a\allowbreak{}s\allowbreak{}t\allowbreak{}(\allowbreak{})}: stacking, placement, tones, actions, dismissal, archive/history review, and shell routing.

This audit verifies the family against real story rendering and runtime behavior. The point is not to invent another notice component. The point is to prove that \texttt{p\allowbreak{}u\allowbreak{}s\allowbreak{}h\allowbreak{}N\allowbreak{}o\allowbreak{}t\allowbreak{}i\allowbreak{}f\allowbreak{}i\allowbreak{}c\allowbreak{}a\allowbreak{}t\allowbreak{}i\allowbreak{}o\allowbreak{}n\allowbreak{}(\allowbreak{})}, \texttt{d\allowbreak{}i\allowbreak{}s\allowbreak{}m\allowbreak{}i\allowbreak{}s\allowbreak{}s\allowbreak{}N\allowbreak{}o\allowbreak{}t\allowbreak{}i\allowbreak{}f\allowbreak{}i\allowbreak{}c\allowbreak{}a\allowbreak{}t\allowbreak{}i\allowbreak{}o\allowbreak{}n\allowbreak{}(\allowbreak{})}, \texttt{t\allowbreak{}i\allowbreak{}c\allowbreak{}k\allowbreak{}N\allowbreak{}o\allowbreak{}t\allowbreak{}i\allowbreak{}f\allowbreak{}i\allowbreak{}c\allowbreak{}a\allowbreak{}t\allowbreak{}i\allowbreak{}o\allowbreak{}n\allowbreak{}s\allowbreak{}(\allowbreak{})}, \texttt{r\allowbreak{}e\allowbreak{}n\allowbreak{}d\allowbreak{}e\allowbreak{}r\allowbreak{}N\allowbreak{}o\allowbreak{}t\allowbreak{}i\allowbreak{}f\allowbreak{}i\allowbreak{}c\allowbreak{}a\allowbreak{}t\allowbreak{}i\allowbreak{}o\allowbreak{}n\allowbreak{}S\allowbreak{}t\allowbreak{}a\allowbreak{}c\allowbreak{}k\allowbreak{}(\allowbreak{})}, \texttt{r\allowbreak{}e\allowbreak{}n\allowbreak{}d\allowbreak{}e\allowbreak{}r\allowbreak{}N\allowbreak{}o\allowbreak{}t\allowbreak{}i\allowbreak{}f\allowbreak{}i\allowbreak{}c\allowbreak{}a\allowbreak{}t\allowbreak{}i\allowbreak{}o\allowbreak{}n\allowbreak{}H\allowbreak{}i\allowbreak{}s\allowbreak{}t\allowbreak{}o\allowbreak{}r\allowbreak{}y\allowbreak{}(\allowbreak{})}, and framed runtime notification routing all still describe one coherent system.

\section{Sponsored Users}\label{record-103-sponsored-users}

\begin{itemize}
\tightlist
\item
  TUI app authors deciding when to use a notification system instead of a local toast, alert, log line, drawer, or modal.
\item
  Docs readers switching DOGFOOD between interactive, static, pipe, and accessible profiles to see what survives without animation or z-order.
\item
  Maintainers who need one focused regression covering notification lifecycle, history, story metadata, and frame-owned routing.
\end{itemize}

\section{Hills}\label{record-103-hills}

\begin{enumerate}
\def\labelenumi{\arabic{enumi}.}
\tightlist
\item
  A builder can inspect DOGFOOD and see live stacks, archived history, and frame-routed notifications as one notification-system family.
\item
  A reader can switch notification stories through every canonical output mode without losing tone, action, ordering, recall, or routing facts.
\item
  A maintainer can run one cycle test that proves notification story docs, component-family doctrine, and runtime behavior have not drifted apart.
\end{enumerate}

\section{Playback Questions}\label{record-103-playback-questions}

\begin{itemize}
\tightlist
\item
  Do DOGFOOD stories show both full notification-system and app-owned transient notification usage?
\item
  Do the stories cover live stack, history review, mixed variants, actions, and framed runtime routing?
\item
  Do rich/static previews keep notification chrome, stacking, placement, and shell context visible?
\item
  Do pipe and accessible lowerings preserve chronological event text without pretending overlay geometry still exists?
\item
  Does the runtime still archive dismissed notifications after their exit phase?
\item
  Do framed apps still route runtime issues and page \texttt{n\allowbreak{}o\allowbreak{}t\allowbreak{}i\allowbreak{}f\allowbreak{}y\allowbreak{}(\allowbreak{})} commands through frame-managed notifications?
\end{itemize}

\section{Requirements}\label{record-103-requirements}

\begin{itemize}
\tightlist
\item
  Keep the \texttt{n\allowbreak{}o\allowbreak{}t\allowbreak{}i\allowbreak{}f\allowbreak{}i\allowbreak{}c\allowbreak{}a\allowbreak{}t\allowbreak{}i\allowbreak{}o\allowbreak{}n\allowbreak{}-\allowbreak{}s\allowbreak{}y\allowbreak{}s\allowbreak{}t\allowbreak{}e\allowbreak{}m} and \texttt{t\allowbreak{}r\allowbreak{}a\allowbreak{}n\allowbreak{}s\allowbreak{}i\allowbreak{}e\allowbreak{}n\allowbreak{}t\allowbreak{}-\allowbreak{}a\allowbreak{}p\allowbreak{}p\allowbreak{}-\allowbreak{}n\allowbreak{}o\allowbreak{}t\allowbreak{}i\allowbreak{}f\allowbreak{}i\allowbreak{}c\allowbreak{}a\allowbreak{}t\allowbreak{}i\allowbreak{}o\allowbreak{}n\allowbreak{}s} DOGFOOD story identities.
\item
  Add a frame-routing variant to the notification-system story.
\item
  Render every notification-system story variant in every canonical story profile.
\item
  Exercise the real notification lifecycle: push, tick, dismiss, archive, render stack, and render history.
\item
  Exercise framed runtime routing through \texttt{c\allowbreak{}r\allowbreak{}e\allowbreak{}a\allowbreak{}t\allowbreak{}e\allowbreak{}F\allowbreak{}r\allowbreak{}a\allowbreak{}m\allowbreak{}e\allowbreak{}d\allowbreak{}A\allowbreak{}p\allowbreak{}p\allowbreak{}(\allowbreak{})} and \texttt{n\allowbreak{}o\allowbreak{}t\allowbreak{}i\allowbreak{}f\allowbreak{}y\allowbreak{}(\allowbreak{})}.
\item
  Keep constrained lowerings plain, chronological, and semantic-preserving.
\end{itemize}

\section{Acceptance Criteria}\label{record-103-acceptance-criteria}

\begin{itemize}
\tightlist
\item
  \texttt{t\allowbreak{}e\allowbreak{}s\allowbreak{}t\allowbreak{}s\allowbreak{}/\allowbreak{}c\allowbreak{}y\allowbreak{}c\allowbreak{}l\allowbreak{}e\allowbreak{}s\allowbreak{}/\allowbreak{}D\allowbreak{}F\allowbreak{}-\allowbreak{}0\allowbreak{}6\allowbreak{}2\allowbreak{}/\allowbreak{}n\allowbreak{}o\allowbreak{}t\allowbreak{}i\allowbreak{}f\allowbreak{}i\allowbreak{}c\allowbreak{}a\allowbreak{}t\allowbreak{}i\allowbreak{}o\allowbreak{}n\allowbreak{}-\allowbreak{}s\allowbreak{}y\allowbreak{}s\allowbreak{}t\allowbreak{}e\allowbreak{}m\allowbreak{}-\allowbreak{}f\allowbreak{}a\allowbreak{}m\allowbreak{}i\allowbreak{}l\allowbreak{}y\allowbreak{}-\allowbreak{}a\allowbreak{}u\allowbreak{}d\allowbreak{}i\allowbreak{}t\allowbreak{}.\allowbreak{}t\allowbreak{}e\allowbreak{}s\allowbreak{}t\allowbreak{}.\allowbreak{}t\allowbreak{}s} proves the cycle doc carries the modern playback sections.
\item
  The test locates DOGFOOD stories for \texttt{n\allowbreak{}o\allowbreak{}t\allowbreak{}i\allowbreak{}f\allowbreak{}i\allowbreak{}c\allowbreak{}a\allowbreak{}t\allowbreak{}i\allowbreak{}o\allowbreak{}n\allowbreak{}-\allowbreak{}s\allowbreak{}y\allowbreak{}s\allowbreak{}t\allowbreak{}e\allowbreak{}m} and \texttt{t\allowbreak{}r\allowbreak{}a\allowbreak{}n\allowbreak{}s\allowbreak{}i\allowbreak{}e\allowbreak{}n\allowbreak{}t\allowbreak{}-\allowbreak{}a\allowbreak{}p\allowbreak{}p\allowbreak{}-\allowbreak{}n\allowbreak{}o\allowbreak{}t\allowbreak{}i\allowbreak{}f\allowbreak{}i\allowbreak{}c\allowbreak{}a\allowbreak{}t\allowbreak{}i\allowbreak{}o\allowbreak{}n\allowbreak{}s}.
\item
  Every documented notification variant renders non-empty output in every profile.
\item
  Rich/static output includes structured visual containment; pipe and accessible output do not depend on box drawing.
\item
  Constrained lowerings preserve live-stack events, history review, framed routing facts, and app-owned transient notices.
\item
  Runtime lifecycle assertions prove notifications can be pushed, ticked, dismissed, archived, stacked, and rendered as history.
\item
  Framed app assertions prove runtime issues and page \texttt{n\allowbreak{}o\allowbreak{}t\allowbreak{}i\allowbreak{}f\allowbreak{}y\allowbreak{}(\allowbreak{})} commands route through frame-managed notifications.
\end{itemize}

\section{Implementation Outline}\label{record-103-implementation-outline}

\begin{enumerate}
\def\labelenumi{\arabic{enumi}.}
\tightlist
\item
  Add a framed-routing variant to the notification-system DOGFOOD story.
\item
  Add a DF-062 cycle test that renders notification stories through the shared story protocol.
\item
  Add runtime assertions for direct notification state and framed app routing.
\item
  Move the backlog note into \texttt{d\allowbreak{}o\allowbreak{}c\allowbreak{}s\allowbreak{}/\allowbreak{}d\allowbreak{}e\allowbreak{}s\allowbreak{}i\allowbreak{}g\allowbreak{}n\allowbreak{}/} and update the v6 lane pointer.
\end{enumerate}

\section{Drift Check}\label{record-103-drift-check}

The runtime already had deep notification support: stack rendering, history rendering, animated entry/exit phases, history archive, frame-owned runtime notifications, and page-level \texttt{n\allowbreak{}o\allowbreak{}t\allowbreak{}i\allowbreak{}f\allowbreak{}y\allowbreak{}(\allowbreak{})} routing.

The drift was explanatory. DOGFOOD showed stack and history, but it did not make frame-routed notifications visible as part of the same family. This slice adds that story variant and binds it to runtime assertions.

\section{Playback}\label{record-103-playback}

\begin{itemize}
\tightlist
\item
  RED: the release lane had only a backlog note, and DOGFOOD did not visibly connect frame-managed runtime notifications to the notification-system family.
\item
  GREEN: the new cycle test renders notification-system and transient notification stories across interactive, static, pipe, and accessible profiles.
\item
  Rich/static previews preserve shell context, stacking, placement, and notification chrome.
\item
  Pipe previews preserve chronological event and history text without box drawing.
\item
  Accessible previews linearize tone, action, recall, and frame-routing facts.
\item
  Runtime assertions prove direct notification lifecycle and framed routing behavior.
\end{itemize}

\section{Retrospective}\label{record-103-retrospective}

This cycle keeps notifications as a system, not a grab bag of overlays. The important boundary is that local \texttt{t\allowbreak{}o\allowbreak{}a\allowbreak{}s\allowbreak{}t\allowbreak{}(\allowbreak{})} is for one directly composed event, while the notification system owns lifecycle, recall, routing, and shell-level review.

